\section{Method}
\label{sec:method}

\subsection{Seeds as capture occasions}
We train $K$ sparse autoencoders on the same activations, in the same order, changing only the random seed of the initialisation, which is the setting in which \citet{paulo2025sparse} found that seeds disagree. Each SAE is one \emph{capture occasion}. A latent is \emph{alive} if it fires on at least $10^{-5}$ of held-out tokens, and it is represented by its unit-norm decoder direction. A \emph{species} is a feature, and the data of the survey is the \emph{capture history} matrix $H\in\{0,1\}^{S\times K}$, where $H_{sr}=1$ when run $r$ contains a latent for species $s$ and $S$ is the number of species seen at least once. Species that no run caught are, by construction, absent from $H$. Estimating how many of them there are is the problem.

\paragraph{Identifying individuals.} An ecologist recognises a recaptured fish by its tag. We need a rule that decides when latents of two runs are the same feature.
\begin{itemize}\setlength\itemsep{0pt}
\item \emph{Oracle} identification (synthetic data only): run $r$ catches planted feature $i$ if one of its alive latents has cosine similarity $\ge\tau$ with $f_i$.
\item \emph{Blind} identification (no ground truth): for every pair of runs, two latents are linked if they are mutual nearest neighbours with decoder cosine $\ge\tau$. Connected components of the resulting graph are species.
\end{itemize}
We use $\tau=\MainTau$ throughout and report $\tau\in\{\TauLo,\TauHi\}$ as robustness checks.

\subsection{Estimators}
Let $Q_j$ be the number of species caught on exactly $j$ of the $K$ occasions, so that $S=\sum_{j\ge1}Q_j$. We compare four families.

\emph{Naive counts.} These are the dictionary size $m$, the number of alive latents in one SAE, and the union $S$ over all seeds, which is the strongest simple baseline because it counts every distinct feature any seed found. We also fit the Michaelis--Menten curve $S(t)=at/(b+t)$ to the seed-accumulation (rarefaction) curve and report its asymptote $a$ \citep{raaijmakers1987statistical,colwell2012models}.

\emph{Equal-catchability models.} These are Chapman's bias-corrected Lincoln--Petersen estimator \citep{chapman1951some}, averaged over all pairs of seeds, and the maximum-likelihood estimator for model $M_t$ \citep{darroch1958multiple}. Both assume every feature is equally easy to catch.

\emph{Heterogeneous-catchability models} (model $M_h$ \citep{otis1978statistical}). These are Chao's lower bound $\hat N=S+\tfrac{K-1}{K}\,Q_1^2/(2Q_2)$ \citep{chao1987estimating}, the first- and second-order jackknives \citep{burnham1978estimation}, and a beta-binomial mixture fitted by maximum likelihood to all observed species \citep{dorazio2003mixture}. Rare features fire on few tokens and contribute little to reconstruction loss, so they are caught less often. That is exactly the heterogeneity these models were built for.

\emph{Singleton-free mixtures} (our adaptation). Blind identification has a failure that tagging does not: a latent that no other seed reproduces becomes a species caught once, whether it is a real but rare feature or seed-specific noise. Microbial ecologists meet the same problem when sequencing errors create spurious singletons that inflate Chao-type estimates \citep{chiu2016estimating}. We therefore fit a catchability mixture only to $Q_2,\dots,Q_K$, the species seen at least twice, by maximising the zero- and one-truncated likelihood. The estimate is $\hat N=S_{\ge2}/\hat P(Y\ge2)$, where $Y$ is the number of seeds that catch a feature. The mixing distribution is a mixture of one to three binomial classes \citep{pledger2000unified}, chosen by AIC, with every class's catch probability bounded below by $p_{\min}=1/K$.

The bound defines what is counted: features that $K$ seeds find at least once on average. It is also what makes the estimate finite. Without it the truncated likelihood barely changes when mass is moved toward $p=0$, so a few stray recaptures can send $\hat N$ to infinity \citep{link2003nonidentifiability}. We call this estimator \emph{singleton-free} and report the unbounded beta-binomial and two-class versions as ablations.

\paragraph{Uncertainty.} Confidence intervals are 95\% percentile intervals from \BootN{} Chao-style bootstrap resamples. We pad $H$ with $\hat N-S$ all-zero histories, resample species with replacement, and re-estimate. Population size under arbitrary heterogeneity is not identifiable \citep{link2003nonidentifiability}: features that no SAE of a given size will ever catch leave no trace. The estimators therefore target features an SAE family can reach, and the synthetic experiments measure how close that is to the truth.

\subsection{Synthetic activations with a planted answer}
Activations are $x=\sum_i a_i f_i$, with $N$ random unit features $f_i\in\mathbb R^{\SynD}$. Feature $i$ is active independently with probability $p_i\propto i^{-\SynAlpha}$, scaled so that \SynMeanActive{} features are active on average, and has magnitude $a_i\sim U(0.5,1.5)$. Planted counts are $N\in\{512,1024,2048\}$, so features outnumber dimensions by a factor of 4--16, the superposition regime \citep{elhage2022toy}. We train TopK SAEs \citep{gao2024scaling} with $k=\SynK$ and dictionary sizes from $N/4$ to $4N$ (capped at 4096), $K=\SynSeeds$ seeds each, for \SynSteps{} steps of 4096 samples (three passes over a fixed set of 1.6M samples). Training uses Adam (learning rate $2\times10^{-3}$, linear decay over the last 20\%), unit-norm decoder rows, and the AuxK loss for dead latents. A second draw of features and data repeats the $N=1024$ sweep.

\subsection{Pythia-70M}
We cache the residual stream after block \RealLayer{} of Pythia-70M \citep{biderman2023pythia} on \RealTrainTokens{} training tokens from 128-token contexts of the Pile \citep{gao2020pile}, using documents 500 onward of the public \texttt{pile-10k} sample. Evaluation uses \RealEvalTokens{} tokens from documents 0--499. The first position of each context is dropped, and activations are scaled so that $\mathbb E\|x\|^2=d$. SAEs use $k=\RealK$, $m\in\{1024,2048,4096,8192\}$, two passes over the data, and $K=\RealSeeds$ seeds per size.
